\begin{abstract}
Environmental forecasts can improve on a persistence reference while
preserving less of the observed dynamic variation. We examined whether
making that distinction explicit changes the set of station--model--horizon
cells passing a post-evaluation screening audit in a rolling-origin
benchmark of daily PM$_{10}$ forecasts from 17 Spanish monitoring stations,
five model families, and seven horizons. Across 595 structured cells, Rule A
required positive persistence-relative RMSE skill and a significant
comparison with persistence; Rule B added study-specific heuristic
thresholds for variance retention and P75 exceedance recall. Rule A retained
277 cells, whereas Rule B retained 8, so 269 cells changed status (97.1\% of
Rule-A cells).
A threshold sensitivity analysis found that 91.3--98.2\% of Rule-A cells
changed across the tested neighbourhood. The results show that error and
dynamic fidelity can produce different cell-screening outcomes under these
stated rules; the magnitude is conditional on this pollutant, benchmark, and
validation design. Here, eligibility denotes only passage of the stated
post-evaluation cell screen, not model ranking or operational acceptance.
\end{abstract}
